% F085｜每种证明就地配图；曲线示例与一般代数论证分工明确。
\FReserve{7.5cm}
\begin{proofbox}
\textbf{\textcolor{CProof}{证明一｜运算法则：指数与对数互为反向观察}}
% F085｜真实数学图；代表函数不构成抽象解的唯一性证明。
\begin{center}
\begin{minipage}{.49\linewidth}\centering
\begin{tikzpicture}
\begin{axis}[width=6.1cm,height=4.8cm,axis lines=middle,
 xmin=-1.3,xmax=3.5,ymin=-0.5,ymax=9,
 xtick={0,1,2,3},xticklabels=\empty,ytick={0,2,4,8},xlabel={$t$},ylabel={$f(t)$},
 title={加变乘：$2^t$},title style={font=\small,yshift=3pt},
 tick label style={font=\small},label style={font=\small},
 ylabel style={at={(axis description cs:0,1)},anchor=south east,xshift=-18pt},
 clip=false,grid=major,grid style={black!8},enlargelimits=false,
 every axis plot/.append style={thick},samples=161]
\addplot[CExplanation,very thick,domain=-1:3] {2^x};
\draw[dashed,black!35] (axis cs:1,0) |- (axis cs:0,2);
\draw[dashed,black!35] (axis cs:2,0) |- (axis cs:0,4);
\draw[dashed,black!35] (axis cs:3,0) |- (axis cs:0,8);
\addplot[only marks,mark=*,mark size=2pt,orange!80!black] coordinates {(1,2) (2,4) (3,8)};

\node[font=\small,anchor=north,yshift=-3pt] at (axis cs:1,-0.5) {$1$};
\node[font=\small,anchor=north,yshift=-3pt] at (axis cs:2,-0.5) {$2$};
\node[font=\small,anchor=north,yshift=-3pt] at (axis cs:3,-0.5) {$3$};

\end{axis}
\end{tikzpicture}
\end{minipage}\hfill
\begin{minipage}{.49\linewidth}\centering
\begin{tikzpicture}
\begin{axis}[width=6.1cm,height=4.8cm,axis lines=middle,
 xmin=-0.3,xmax=8.8,ymin=-2,ymax=3.7,
 xtick={1,2,4,8},xticklabels=\empty,ytick={-1,0,1,2,3},xlabel={$t$},ylabel={$f(t)$},
 title={乘变加：$\log_2t$},title style={font=\small,yshift=3pt},
 tick label style={font=\small},label style={font=\small},
 ylabel style={at={(axis description cs:0,1)},anchor=south east,xshift=-18pt},
 clip=false,grid=major,grid style={black!8},enlargelimits=false,
 every axis plot/.append style={thick},samples=161]
\addplot[CExplanation,very thick,domain=.3:8] {ln(x)/ln(2)};
\draw[dashed,black!35] (axis cs:2,0) |- (axis cs:0,1);
\draw[dashed,black!35] (axis cs:4,0) |- (axis cs:0,2);
\draw[dashed,black!35] (axis cs:8,0) |- (axis cs:0,3);
\addplot[only marks,mark=*,mark size=2pt,orange!80!black] coordinates {(2,1) (4,2) (8,3)};

\node[font=\small,anchor=north,yshift=-3pt] at (axis cs:1,-2) {$1$};
\node[font=\small,anchor=north,yshift=-3pt] at (axis cs:2,-2) {$2$};
\node[font=\small,anchor=north,yshift=-3pt] at (axis cs:4,-2) {$4$};
\node[font=\small,anchor=north,yshift=-3pt] at (axis cs:8,-2) {$8$};

\end{axis}
\end{tikzpicture}
\end{minipage}
\par\small 同一组数反向阅读：$2^{1+2}=2\cdot4=8$；$\log_2(2\cdot4)=1+2=3$。
\end{center}

左图把输入相加变成高度相乘，右图反向读取相同的数。这一观察对应下列对所有合法输入成立的运算法则：
\[
\begin{aligned}
f(x)=a^x&:\quad f(x+y)=a^{x+y}=a^xa^y=f(x)f(y),\\
f(x)=\log_a x&:\quad f(xy)=\log_a x+\log_a y=f(x)+f(y).
\end{aligned}
\]
指数取 $a>0$；对数取 $x,y>0,a>0,a\ne1$。图用 $a=2$ 展示，证明使用任意合法底数。
\end{proofbox}

\FReserve{7.5cm}
\begin{proofbox}
\textbf{\textcolor{CProof}{证明一（续）｜幂与正比例分别保留乘法、加法}}
% F085｜真实数学图；代表函数不构成抽象解的唯一性证明。
\begin{center}
\begin{minipage}{.49\linewidth}\centering
\begin{tikzpicture}
\begin{axis}[width=6.1cm,height=4.8cm,axis lines=middle,
 xmin=-0.2,xmax=3.5,ymin=-0.5,ymax=10.8,
 xtick={0,1,2,3},xticklabels=\empty,ytick={0,1,4,9},xlabel={$t$},ylabel={$f(t)$},
 title={幂：乘法作用于横坐标与高度},title style={font=\small,yshift=3pt},
 tick label style={font=\small},label style={font=\small},
 ylabel style={at={(axis description cs:0,1)},anchor=south east,xshift=-18pt},
 clip=false,grid=major,grid style={black!8},enlargelimits=false,
 every axis plot/.append style={thick},samples=161]
\addplot[CExplanation,very thick,domain=0:3.1] {x^2};
\draw[dashed,black!35] (axis cs:1,0) |- (axis cs:0,1);
\draw[dashed,black!35] (axis cs:2,0) |- (axis cs:0,4);
\draw[dashed,black!35] (axis cs:3,0) |- (axis cs:0,9);
\addplot[only marks,mark=*,mark size=2pt,orange!80!black] coordinates {(1,1) (2,4) (3,9)};

\node[font=\small,anchor=north,yshift=-3pt] at (axis cs:1,-0.5) {$1$};
\node[font=\small,anchor=north,yshift=-3pt] at (axis cs:2,-0.5) {$2$};
\node[font=\small,anchor=north,yshift=-3pt] at (axis cs:3,-0.5) {$3$};

\end{axis}
\end{tikzpicture}
\end{minipage}\hfill
\begin{minipage}{.49\linewidth}\centering
\begin{tikzpicture}
\begin{axis}[width=6.1cm,height=4.8cm,axis lines=middle,
 xmin=-0.2,xmax=3.5,ymin=-0.4,ymax=7.5,
 xtick={0,1,2,3},xticklabels=\empty,ytick={0,2,4,6},xlabel={$t$},ylabel={$f(t)$},
 title={正比例：加法作用于横坐标与高度},title style={font=\small,yshift=3pt},
 tick label style={font=\small},label style={font=\small},
 ylabel style={at={(axis description cs:0,1)},anchor=south east,xshift=-18pt},
 clip=false,grid=major,grid style={black!8},enlargelimits=false,
 every axis plot/.append style={thick},samples=161]
\addplot[CExplanation,very thick,domain=0:3.1] {2*x};
\draw[dashed,black!35] (axis cs:1,0) |- (axis cs:0,2);
\draw[dashed,black!35] (axis cs:2,0) |- (axis cs:0,4);
\draw[dashed,black!35] (axis cs:3,0) |- (axis cs:0,6);
\addplot[only marks,mark=*,mark size=2pt,orange!80!black] coordinates {(1,2) (2,4) (3,6)};

\node[font=\small,anchor=north,yshift=-3pt] at (axis cs:1,-0.4) {$1$};
\node[font=\small,anchor=north,yshift=-3pt] at (axis cs:2,-0.4) {$2$};
\node[font=\small,anchor=north,yshift=-3pt] at (axis cs:3,-0.4) {$3$};

\end{axis}
\end{tikzpicture}
\end{minipage}
\par\small 左图 $f(\sqrt2\cdot\sqrt2)=f(\sqrt2)^2=4$；右图 $f(1+2)=f(1)+f(2)=6$。
\end{center}

\[
\begin{aligned}
f(x)=x^a&:\quad f(xy)=(xy)^a=x^ay^a=f(x)f(y)\quad(x,y>0),\\
f(x)=kx&:\quad f(x+y)=k(x+y)=kx+ky=f(x)+f(y).
\end{aligned}
\]
这些证明只说明典型函数满足模型；反过来确定唯一解析式，还须题目提供充分条件。
\end{proofbox}

\FReserve{7.5cm}
\begin{proofbox}
\textbf{\textcolor{CProof}{证明二｜余弦和差相加：弦中点的横坐标}}
% F085｜真实数学图；代表函数不构成抽象解的唯一性证明。
\begin{center}
\begin{tikzpicture}[x=2.6cm,y=2.6cm,>=stealth]
\draw[->,black!65] (-.3,0)--(1.5,0) node[right,font=\small] {横轴};
\draw[->,black!65] (0,-.15)--(0,1.3);
\draw[CExplanation,thick] (0,0)--(30:1) node[right,font=\small,xshift=4pt] {$P_-$};
\draw[CExplanation,thick] (0,0)--(60:1) node[above right,font=\small] {$M$};
\draw[CExplanation,thick] (0,0)--(90:1) node[above right,font=\small] {$P_+$};
\draw[CExplanation,thick] (0:1) arc (0:90:1);
\draw[orange!80!black,very thick] (30:1)--(90:1);
\fill[orange!80!black] (.4330127,.75) circle (.025);
\draw[orange!80!black] (.47,.78)--(1.1,1.15) node[right,font=\small] {$Q$（弦中点）};
\draw[dashed] (30:1)--(.8660254,-.25);
\draw[dashed] (90:1)--(0,-.25);
\draw[dashed] (.4330127,.75)--(.4330127,-.25);
\draw[<->,thick,orange!80!black] (0,-.18)--(.8660254,-.18);
\node[below,font=\small] at (.433,-.25) {两端横坐标的平均值是 $Q$ 的横坐标};
\draw (0.23,0) arc (0:60:.23);
\node[font=\small] at (.25,.26) {$x$};
\node[font=\small,below left] at (0,0) {$O$};
\node[font=\small,align=left,anchor=west] at (1.9,.58) {$P_\pm=(\cos(x\pm y),\sin(x\pm y))$\\[5pt]
$OM=1,\quad OQ=\cos y$\\[5pt]
$Q$ 在方向角 $x$ 的射线上\\[5pt]
$Q$ 的横坐标 $=\cos x\cos y$};
\end{tikzpicture}
\par\small 图取 $x=\pi/3,y=\pi/6$。单位圆的弦中点给出 $[\cos(x+y)+\cos(x-y)]/2=\cos x\cos y$。
\end{center}

图中 $P_+,P_-$ 的横坐标平均值是弦中点 $Q$ 的横坐标；射线 $OQ$ 的方向角为 $x$，$OQ=\cos y$。于是看见两侧投影的平均值 $\cos x\cos y$。

\textbf{任意角的代数证明：}分别展开
\[
\begin{aligned}
\cos(x+y)&=\cos x\cos y-\sin x\sin y,\\
\cos(x-y)&=\cos x\cos y+\sin x\sin y.
\end{aligned}
\]
相加，异号的 $\sin x\sin y$ 抵消，得到
\[\boxed{\cos(x+y)+\cos(x-y)=2\cos x\cos y}.\]
\end{proofbox}

\FReserve{7.5cm}
\begin{proofbox}
\textbf{\textcolor{CProof}{证明二（续）｜正弦和差相乘：平方差的面积拆拼}}
% F085｜真实数学图；代表函数不构成抽象解的唯一性证明。
\begin{center}
\begin{tikzpicture}[x=1.3cm,y=1.3cm,>=stealth]
\fill[CExplanation!18] (0,0) rectangle (3,3);
\fill[white] (0,0) rectangle (1,1);
\draw[CExplanation,thick] (0,0) rectangle (3,3);
\draw[orange!80!black,thick] (0,0) rectangle (1,1);
\draw[dashed] (1,0)--(1,3);
\node[font=\small] at (2,1.5) {$A(A-B)$};
\node[font=\small,rotate=90] at (.5,2) {$B(A-B)$};
\node[font=\small] at (.5,.5) {$B^2$};
\draw[<->] (0,3.25)--node[above,font=\small] {$A$}(3,3.25);
\draw[<->] (-.25,0)--node[left,font=\small] {$B$}(-.25,1);
\draw[->,thick] (3.4,1.5)--(4.3,1.5);
\fill[CExplanation!18] (4.7,0) rectangle (8.7,2);
\draw[CExplanation,thick] (4.7,0) rectangle (8.7,2);
\draw[dashed] (7.7,0)--(7.7,2);
\node[font=\small] at (6.2,1) {$A(A-B)$};
\node[font=\small,rotate=90] at (8.2,1) {$B(A-B)$};
\draw[<->] (4.7,-.25)--node[below,font=\small] {$A+B$}(8.7,-.25);
\draw[<->] (8.95,0)--node[right,font=\small] {$A-B$}(8.95,2);
\end{tikzpicture}
\par\small 面积拆拼：$A^2-B^2=(A+B)(A-B)$。这里示意取 $A=3,B=1$；正弦证明中取 $A=\sin x\cos y,B=\cos x\sin y$，任意符号情形由代数恒等式保证。
\end{center}

设 $A=\sin x\cos y,B=\cos x\sin y$。由和差公式，两个正弦值就是 $A+B$ 与 $A-B$。图中两块剩余面积拼成长方形，对应
\[
\begin{aligned}
\sin(x+y)\sin(x-y)&=(A+B)(A-B)=A^2-B^2\\
&=\sin^2x(1-\sin^2y)-(1-\sin^2x)\sin^2y\\
&=\boxed{\sin^2x-\sin^2y}.
\end{aligned}
\]
面积图展示 $A>B>0$ 的情形；上面的代数恒等式不受 $A,B$ 正负限制。
\end{proofbox}

\FReserve{7.5cm}
\begin{proofbox}
\textbf{\textcolor{CProof}{证明三｜正切与余切：同一射线的两种比值}}
% F085｜真实数学图；代表函数不构成抽象解的唯一性证明。
\begin{center}
\begin{tikzpicture}[x=2.3cm,y=2.3cm,>=stealth]
\draw[->] (-.1,0)--(2.1,0) node[right,font=\small] {$u$};
\draw[->] (0,-.1)--(0,1.8) node[above,font=\small] {$v$};
\draw[CExplanation,thick] (1,0) arc(0:90:1);
\draw[black!40] (1,0)--(1,1.7);
\draw[black!40] (0,1)--(1.9,1);
\draw[CExplanation,very thick] (0,0)--(1.7,1.7);
\fill[CExplanation] (.7071068,.7071068) circle(.022) node[above left,font=\small] {$P$};
\draw[dashed] (.7071068,0)--(.7071068,.7071068)--(0,.7071068);
\draw[orange!80!black,very thick] (1,0)--(1,1);
\draw[green!45!black,very thick] (0,1)--(1,1);
\draw (.25,0) arc(0:45:.25);
\node[font=\small] at (.38,.13) {$\theta$};
\node[font=\small,anchor=west,text=orange!80!black] at (1.07,.5) {$\tan\theta$};
\node[font=\small,above,text=green!45!black] at (.5,1.03) {$\cot\theta$};
\node[font=\small,anchor=west,align=left] at (2.6,1.2) {$P=(\cos\theta,\sin\theta)$\\[7pt]
竖直切线 $u=1$：\\[3pt]$\tan\theta=\sin\theta/\cos\theta$\\[7pt]
水平切线 $v=1$：\\[3pt]$\cot\theta=\cos\theta/\sin\theta$};
\end{tikzpicture}
\par\small 取 $\theta=x+y$。同一条射线的纵横比与横纵比互为倒数；相应分母为 $0$ 时，比值无定义。图示 $\theta=\pi/4$。
\end{center}

对 $\theta=x+y$，先分别写成纵横比与横纵比，再代入和角公式：
\[
\begin{aligned}
\tan(x+y)&=\frac{\sin x\cos y+\cos x\sin y}{\cos x\cos y-\sin x\sin y}
=\frac{\tan x+\tan y}{1-\tan x\tan y},\\[5pt]
\cot(x+y)&=\frac{\cos x\cos y-\sin x\sin y}{\sin x\cos y+\cos x\sin y}
=\frac{\cot x\cot y-1}{\cot x+\cot y}.
\end{aligned}
\]
正切一式同除以 $\cos x\cos y$；余切一式同除以 $\sin x\sin y$。相应的 $x,y,x+y$ 处函数值均须有定义，所有除数均须非零。
\end{proofbox}

\FReserve{7.5cm}
\begin{proofbox}
\textbf{\textcolor{CProof}{证明四｜特殊值代入：从模型直接推出对称性}}
% F085｜真实数学图；代表函数不构成抽象解的唯一性证明。
\begin{center}
\begin{minipage}{.49\linewidth}\centering
\begin{tikzpicture}
\begin{axis}[width=6.1cm,height=4.8cm,axis lines=middle,
 xmin=-3.7,xmax=3.7,ymin=-1.4,ymax=1.6,
 xtick={-3,0,3},xticklabels=\empty,ytick={-1,0,1},xlabel={$t$},ylabel={$f(t)$},
 title={余弦代表：关于纵轴对称},title style={font=\small,yshift=3pt},
 tick label style={font=\small},label style={font=\small},
 ylabel style={at={(axis description cs:0,1)},anchor=south east,xshift=-18pt},
 clip=false,grid=major,grid style={black!8},enlargelimits=false,
 every axis plot/.append style={thick},samples=161]
\addplot[CExplanation,very thick,domain=-3.14:3.14] {cos(deg(x))};
\draw[dashed,black!35] (axis cs:-1.0471975511965976,0) |- (axis cs:0,0.5);
\draw[dashed,black!35] (axis cs:1.0471975511965976,0) |- (axis cs:0,0.5);
\addplot[only marks,mark=*,mark size=2pt,orange!80!black] coordinates {(-1.04719755,0.5) (1.04719755,0.5)};

\node[font=\small,anchor=north,yshift=-3pt] at (axis cs:-3,-1.4) {$-3$};
\node[font=\small,anchor=north,yshift=-3pt] at (axis cs:3,-1.4) {$3$};

\end{axis}
\end{tikzpicture}
\end{minipage}\hfill
\begin{minipage}{.49\linewidth}\centering
\begin{tikzpicture}
\begin{axis}[width=6.1cm,height=4.8cm,axis lines=middle,
 xmin=-3.7,xmax=3.7,ymin=-1.5,ymax=1.5,
 xtick={-3,0,3},xticklabels=\empty,ytick={-1,0,1},xlabel={$t$},ylabel={$f(t)$},
 title={正弦代表：关于原点对称},title style={font=\small,yshift=3pt},
 tick label style={font=\small},label style={font=\small},
 ylabel style={at={(axis description cs:0,1)},anchor=south east,xshift=-18pt},
 clip=false,grid=major,grid style={black!8},enlargelimits=false,
 every axis plot/.append style={thick},samples=161]
\addplot[CExplanation,very thick,domain=-3.14:3.14] {sin(deg(x))};
\draw[dashed,black!35] (axis cs:-0.5235987755982988,0) |- (axis cs:0,-0.5);
\draw[dashed,black!35] (axis cs:0.5235987755982988,0) |- (axis cs:0,0.5);
\addplot[only marks,mark=*,mark size=2pt,orange!80!black] coordinates {(-0.523598776,-0.5) (0.523598776,0.5)};

\node[font=\small,anchor=north,yshift=-3pt] at (axis cs:-3,-1.5) {$-3$};
\node[font=\small,anchor=north,yshift=-3pt] at (axis cs:3,-1.5) {$3$};

\end{axis}
\end{tikzpicture}
\end{minipage}
\par\small 特殊值代入得到的是对称关系。这里用代表函数展示，证明对抽象模型直接进行。
\end{center}

\textbf{指数型（实数域、不恒为零）：}令 $y=0$，从 $f(x)=f(x)f(0)$ 选取一个非零函数值得 $f(0)=1$。令 $y=-x$，得 $f(x)f(-x)=1$。又 $f(x)=f(x/2)^2>0$。

\textbf{对数型（正数域）：}令 $y=1$ 得 $f(1)=0$；令 $y=1/x$ 得 $f(1/x)=-f(x)$。

\textbf{幂型（正数域、不恒为零）：}令 $y=1$ 得 $f(1)=1$；令 $y=1/x$ 得 $f(x)f(1/x)=1$，且 $f(x)=f(\sqrt x)^2>0$。

\textbf{正比例型（实数域）：}令 $y=0$ 得 $f(0)=0$；令 $y=-x$ 得 $f(-x)=-f(x)$。

\textbf{余弦型（实数域、不恒为零）：}令 $y=0$ 得 $f(0)=1$；令 $x=0$ 得 $f(-y)=f(y)$，对应左图的轴对称。

\textbf{正弦型（实数域）：}令 $y=0$ 得 $f(0)^2=0$。令 $x=0,y=\pm t$，分别得
\[f(t)f(-t)=-f(t)^2=-f(-t)^2.\]
所以 $[f(t)+f(-t)]^2=0$，即 $f(-t)=-f(t)$，对应右图的中心对称。

以上结论由方程推出，\textbf{不需要预先假设连续或猜出解析式}；图中曲线只帮助看见关系。
\end{proofbox}

\FReserve{7.5cm}
\begin{proofbox}
\textbf{\textcolor{CProof}{重要推论 1｜相加输入：倍数与相反数}}
% F085｜真实数学图；代表函数不构成抽象解的唯一性证明。
\begin{center}
\begin{minipage}{.49\linewidth}\centering
\begin{tikzpicture}
\begin{axis}[width=6.1cm,height=4.8cm,axis lines=middle,
 xmin=-2.5,xmax=2.6,ymin=-0.3,ymax=4.8,
 xtick={-2,-1,0,1,2},xticklabels=\empty,ytick={0,1,2,4},xlabel={$t$},ylabel={$f(t)$},
 title={指数：倍数变乘方},title style={font=\small,yshift=3pt},
 tick label style={font=\small},label style={font=\small},
 ylabel style={at={(axis description cs:0,1)},anchor=south east,xshift=-18pt},
 clip=false,grid=major,grid style={black!8},enlargelimits=false,
 every axis plot/.append style={thick},samples=161]
\addplot[CExplanation,very thick,domain=-2:2.1] {2^x};
\draw[dashed,black!35] (axis cs:-1,0) |- (axis cs:0,0.5);
\draw[dashed,black!35] (axis cs:1,0) |- (axis cs:0,2);
\draw[dashed,black!35] (axis cs:2,0) |- (axis cs:0,4);
\addplot[only marks,mark=*,mark size=2pt,orange!80!black] coordinates {(-1,0.5) (1,2) (2,4)};

\node[font=\small,anchor=north,yshift=-3pt] at (axis cs:-2,-0.3) {$-2$};
\node[font=\small,anchor=north,yshift=-3pt] at (axis cs:-1,-0.3) {$-1$};
\node[font=\small,anchor=north,yshift=-3pt] at (axis cs:1,-0.3) {$1$};
\node[font=\small,anchor=north,yshift=-3pt] at (axis cs:2,-0.3) {$2$};

\end{axis}
\end{tikzpicture}
\end{minipage}\hfill
\begin{minipage}{.49\linewidth}\centering
\begin{tikzpicture}
\begin{axis}[width=6.1cm,height=4.8cm,axis lines=middle,
 xmin=-2.5,xmax=2.6,ymin=-4.8,ymax=4.8,
 xtick={-2,-1,0,1,2},xticklabels=\empty,ytick={-4,-2,0,2,4},xlabel={$t$},ylabel={$f(t)$},
 title={正比例：倍数仍是倍数},title style={font=\small,yshift=3pt},
 tick label style={font=\small},label style={font=\small},
 ylabel style={at={(axis description cs:0,1)},anchor=south east,xshift=-18pt},
 clip=false,grid=major,grid style={black!8},enlargelimits=false,
 every axis plot/.append style={thick},samples=161]
\addplot[CExplanation,very thick,domain=-2:2.1] {2*x};
\draw[dashed,black!35] (axis cs:-1,0) |- (axis cs:0,-2);
\draw[dashed,black!35] (axis cs:1,0) |- (axis cs:0,2);
\draw[dashed,black!35] (axis cs:2,0) |- (axis cs:0,4);
\addplot[only marks,mark=*,mark size=2pt,orange!80!black] coordinates {(-1,-2) (1,2) (2,4)};

\node[font=\small,anchor=north,yshift=-3pt] at (axis cs:-2,-4.8) {$-2$};
\node[font=\small,anchor=north,yshift=-3pt] at (axis cs:-1,-4.8) {$-1$};
\node[font=\small,anchor=north,yshift=-3pt] at (axis cs:1,-4.8) {$1$};
\node[font=\small,anchor=north,yshift=-3pt] at (axis cs:2,-4.8) {$2$};

\end{axis}
\end{tikzpicture}
\end{minipage}
\par\small 同取 $f(1)=2$：指数代表有 $f(-1)=1/2$；正比例代表有 $f(-1)=-2$。相反输入的高度关系不同。
\end{center}

反复使用加法型关系，再用相反数关系扩展到负整数，得到
\[
\begin{array}{lll}
\text{指数型：}&f(nx)=[f(x)]^n,&f(-x)=1/f(x);\\[4pt]
\text{正比例型：}&f(nx)=nf(x),&f(-x)=-f(x).
\end{array}
\]
这里 $n\in\mathbb Z$，关系在全体实数上成立；指数型须不恒为零，因此 $f(x)>0$，$n=0$ 和负整数幂均有意义。
\end{proofbox}

\FReserve{7.5cm}
\begin{proofbox}
\textbf{\textcolor{CProof}{重要推论 2｜相乘输入：乘方与倒数}}
% F085｜真实数学图；代表函数不构成抽象解的唯一性证明。
\begin{center}
\begin{minipage}{.49\linewidth}\centering
\begin{tikzpicture}
\begin{axis}[width=6.1cm,height=4.8cm,axis lines=middle,
 xmin=-0.2,xmax=4.6,ymin=-2,ymax=3.1,
 xtick={1,2,4},xticklabels=\empty,ytick={-1,0,1,2},xlabel={$t$},ylabel={$f(t)$},
 title={对数：$f(t)=\log_2t$},title style={font=\small,yshift=3pt},
 tick label style={font=\small},label style={font=\small},
 ylabel style={at={(axis description cs:0,1)},anchor=south east,xshift=-18pt},
 clip=false,grid=major,grid style={black!8},enlargelimits=false,
 every axis plot/.append style={thick},samples=161]
\addplot[CExplanation,very thick,domain=.3:4.1] {ln(x)/ln(2)};
\draw[dashed,black!35] (axis cs:0.5,0) |- (axis cs:0,-1);
\draw[dashed,black!35] (axis cs:1,0) |- (axis cs:0,0);
\draw[dashed,black!35] (axis cs:2,0) |- (axis cs:0,1);
\draw[dashed,black!35] (axis cs:4,0) |- (axis cs:0,2);
\addplot[only marks,mark=*,mark size=2pt,orange!80!black] coordinates {(0.5,-1) (1,0) (2,1) (4,2)};

\node[font=\small,anchor=north,yshift=-3pt] at (axis cs:1,-2) {$1$};
\node[font=\small,anchor=north,yshift=-3pt] at (axis cs:2,-2) {$2$};
\node[font=\small,anchor=north,yshift=-3pt] at (axis cs:4,-2) {$4$};

\end{axis}
\end{tikzpicture}
\end{minipage}\hfill
\begin{minipage}{.49\linewidth}\centering
\begin{tikzpicture}
\begin{axis}[width=6.1cm,height=4.8cm,axis lines=middle,
 xmin=-0.2,xmax=4.6,ymin=-1,ymax=18,
 xtick={1,2,4},xticklabels=\empty,ytick={0,4,8,16},xlabel={$t$},ylabel={$f(t)$},
 title={幂型：$f(t)=t^2$},title style={font=\small,yshift=3pt},
 tick label style={font=\small},label style={font=\small},
 ylabel style={at={(axis description cs:0,1)},anchor=south east,xshift=-18pt},
 clip=false,grid=major,grid style={black!8},enlargelimits=false,
 every axis plot/.append style={thick},samples=161]
\addplot[CExplanation,very thick,domain=.2:4.1] {x^2};
\draw[dashed,black!35] (axis cs:0.5,0) |- (axis cs:0,0.25);
\draw[dashed,black!35] (axis cs:1,0) |- (axis cs:0,1);
\draw[dashed,black!35] (axis cs:2,0) |- (axis cs:0,4);
\draw[dashed,black!35] (axis cs:4,0) |- (axis cs:0,16);
\addplot[only marks,mark=*,mark size=2pt,orange!80!black] coordinates {(0.5,0.25) (1,1) (2,4) (4,16)};

\node[font=\small,anchor=north,yshift=-3pt] at (axis cs:1,-1) {$1$};
\node[font=\small,anchor=north,yshift=-3pt] at (axis cs:2,-1) {$2$};
\node[font=\small,anchor=north,yshift=-3pt] at (axis cs:4,-1) {$4$};

\end{axis}
\end{tikzpicture}
\end{minipage}
\par\small 看 $t=2\to4=2^2$：左图高度加倍 $1\to2$，右图高度平方 $4\to16$；倒数输入 $1/2$ 分别对应负值与倒数值。
\end{center}

把 $x^n$ 拆成相同因子的乘积，$n=0$ 用 $f(1)$，负整数再用倒数关系：
\[
\begin{array}{lll}
\text{对数型：}&f(x^n)=nf(x),&f(1/x)=-f(x);\\[4pt]
\text{幂函数型：}&f(x^n)=[f(x)]^n,&f(1/x)=1/f(x).
\end{array}
\]
这里 $x>0,n\in\mathbb Z$；幂型须不恒为零。左图是高度乘 $n$，右图是高度的 $n$ 次幂，不能混用。
\end{proofbox}

\FReserve{7.5cm}
\begin{proofbox}
\textbf{\textcolor{CProof}{重要推论 3｜余弦的二倍输入，正弦的中点半差}}
% F085｜真实数学图；代表函数不构成抽象解的唯一性证明。
\begin{center}
\begin{minipage}{.49\linewidth}\centering
\begin{tikzpicture}
\begin{axis}[width=6.1cm,height=4.8cm,axis lines=middle,
 xmin=-0.2,xmax=3.6,ymin=-1.4,ymax=1.5,
 xtick={0,1.047198,2.094395,3.141593},xticklabels=\empty,ytick={-1,0,1},xlabel={$t$},ylabel={$f(t)$},
 title={余弦：输入加倍，高度平方后变换},title style={font=\small,yshift=3pt},
 tick label style={font=\small},label style={font=\small},
 ylabel style={at={(axis description cs:0,1)},anchor=south east,xshift=-18pt},
 clip=false,grid=major,grid style={black!8},enlargelimits=false,
 every axis plot/.append style={thick},samples=161]
\addplot[CExplanation,very thick,domain=0:3.2] {cos(deg(x))};
\draw[dashed,black!35] (axis cs:1.0471975511965976,0) |- (axis cs:0,0.5);
\draw[dashed,black!35] (axis cs:2.0943951023931953,0) |- (axis cs:0,-0.5);
\addplot[only marks,mark=*,mark size=2pt,orange!80!black] coordinates {(1.04719755,0.5) (2.0943951,-0.5)};

\node[font=\small,anchor=north,yshift=-3pt] at (axis cs:1.047198,-1.4) {$\pi/3$};
\node[font=\small,anchor=north,yshift=-3pt] at (axis cs:2.094395,-1.4) {$2\pi/3$};
\node[font=\small,anchor=north,yshift=-3pt] at (axis cs:3.141593,-1.4) {$\pi$};

\end{axis}
\end{tikzpicture}
\end{minipage}\hfill
\begin{minipage}{.49\linewidth}\centering
\begin{tikzpicture}
\begin{axis}[width=6.1cm,height=4.8cm,axis lines=middle,
 xmin=-0.2,xmax=2.9,ymin=-0.3,ymax=1.5,
 xtick={0,0.523599,1.570796},xticklabels=\empty,ytick={0,0.5,1},xlabel={$t$},ylabel={$f(t)$},
 title={正弦：用中点和半差配对},title style={font=\small,yshift=3pt},
 tick label style={font=\small},label style={font=\small},
 ylabel style={at={(axis description cs:0,1)},anchor=south east,xshift=-18pt},
 clip=false,grid=major,grid style={black!8},enlargelimits=false,
 every axis plot/.append style={thick},samples=161]
\addplot[CExplanation,very thick,domain=0:2.5] {sin(deg(x))};
\draw[dashed,black!35] (axis cs:0.5235987755982988,0) |- (axis cs:0,0.5);
\draw[dashed,black!35] (axis cs:1.5707963267948966,0) |- (axis cs:0,1);
\addplot[only marks,mark=*,mark size=2pt,orange!80!black] coordinates {(0.523598776,0.5) (1.57079633,1)};

\node[font=\small,anchor=north,yshift=-3pt] at (axis cs:0.523599,-0.3) {$\pi/6$};
\node[font=\small,anchor=north,yshift=-3pt] at (axis cs:1.570796,-0.3) {$\pi/2$};

\end{axis}
\end{tikzpicture}
\end{minipage}
\par\small 左图 $t=\pi/3\to2\pi/3$；右图 $m=\pi/2,n=\pi/6$，中点 $x=\pi/3$，半差 $y=\pi/6$，故乘积 $1/2=3/4-1/4$。
\end{center}

\textbf{余弦型：}不恒为零时 $f(0)=1$。令 $y=x$ 得 $f(2x)+1=2f(x)^2$，即
\[\boxed{f(2x)=2[f(x)]^2-1}.\]
\textbf{正弦型：}若要计算 $f(m)f(n)$，在数轴上取中点 $x=(m+n)/2$、半差 $y=(m-n)/2$，便有 $x+y=m,x-y=n$，所以
\[\boxed{f(m)f(n)=f^2\!\left(\frac{m+n}{2}\right)-f^2\!\left(\frac{m-n}{2}\right)}.\]
两者均按实数域上的模型讨论；抽象函数的推论并不要求其图像就是图中的三角曲线。
\end{proofbox}

\FReserve{7.5cm}
\begin{proofbox}
\textbf{\textcolor{CProof}{重要推论 4｜正切、余切的二倍输入与无定义边界}}
% F085｜真实数学图；代表函数不构成抽象解的唯一性证明。
\begin{center}
\begin{minipage}{.49\linewidth}\centering
\begin{tikzpicture}
\begin{axis}[width=6.1cm,height=4.8cm,axis lines=middle,
 xmin=-0.1,xmax=1.5,ymin=-0.2,ymax=3.8,
 xtick={0,0.523599,1.047198},xticklabels=\empty,ytick={0,1,2,3},xlabel={$t$},ylabel={$f(t)$},
 title={正切二倍角：$\tan t$},title style={font=\small,yshift=3pt},
 tick label style={font=\small},label style={font=\small},
 ylabel style={at={(axis description cs:0,1)},anchor=south east,xshift=-18pt},
 clip=false,grid=major,grid style={black!8},enlargelimits=false,
 every axis plot/.append style={thick},samples=161]
\addplot[CExplanation,very thick,domain=0:1.27] {tan(deg(x))};
\draw[dashed,black!35] (axis cs:0.5235987755982988,0) |- (axis cs:0,0.5773502691896258);
\draw[dashed,black!35] (axis cs:1.0471975511965976,0) |- (axis cs:0,1.7320508075688772);
\addplot[only marks,mark=*,mark size=2pt,orange!80!black] coordinates {(0.523598776,0.577350269) (1.04719755,1.73205081)};

\node[font=\small,anchor=north,yshift=-3pt] at (axis cs:0.523599,-0.2) {$\pi/6$};
\node[font=\small,anchor=north,yshift=-3pt] at (axis cs:1.047198,-0.2) {$\pi/3$};

\end{axis}
\end{tikzpicture}
\end{minipage}\hfill
\begin{minipage}{.49\linewidth}\centering
\begin{tikzpicture}
\begin{axis}[width=6.1cm,height=4.8cm,axis lines=middle,
 xmin=-0.1,xmax=1.5,ymin=-0.2,ymax=3.8,
 xtick={0,0.523599,1.047198},xticklabels=\empty,ytick={0,1,2,3},xlabel={$t$},ylabel={$f(t)$},
 title={余切二倍角：$\cot t$},title style={font=\small,yshift=3pt},
 tick label style={font=\small},label style={font=\small},
 ylabel style={at={(axis description cs:0,1)},anchor=south east,xshift=-18pt},
 clip=false,grid=major,grid style={black!8},enlargelimits=false,
 every axis plot/.append style={thick},samples=161]
\addplot[CExplanation,very thick,domain=.27:1.45] {cos(deg(x))/sin(deg(x))};
\draw[dashed,black!35] (axis cs:0.5235987755982988,0) |- (axis cs:0,1.7320508075688772);
\draw[dashed,black!35] (axis cs:1.0471975511965976,0) |- (axis cs:0,0.5773502691896258);
\addplot[only marks,mark=*,mark size=2pt,orange!80!black] coordinates {(0.523598776,1.73205081) (1.04719755,0.577350269)};

\node[font=\small,anchor=north,yshift=-3pt] at (axis cs:0.523599,-0.2) {$\pi/6$};
\node[font=\small,anchor=north,yshift=-3pt] at (axis cs:1.047198,-0.2) {$\pi/3$};

\end{axis}
\end{tikzpicture}
\end{minipage}
\par\small 同一输入 $\pi/6\to\pi/3$：正切高度 $1/\sqrt3\to\sqrt3$，余切高度 $\sqrt3\to1/\sqrt3$；二倍输入处必须有定义。
\end{center}

在原关系中取 $y=x$，直接得到
\[
\text{正切型： }\boxed{f(2x)=\frac{2f(x)}{1-[f(x)]^2}},\qquad
\text{余切型： }\boxed{f(2x)=\frac{[f(x)]^2-1}{2f(x)}}.
\]
前者要求 $f(x)^2\ne1$；后者要求 $f(x)\ne0$，并且两者均须在 $x,2x$ 处有定义。比如 $\tan(\pi/4)=1$，但 $\tan(\pi/2)$ 不存在，不能给它算出一个“函数值”。
\end{proofbox}

